\section{Mechanism Taxonomy and Genealogy}
\label{sec:taxonomy}

\subsection{From algorithm names to computational mechanisms}

Optimization taxonomies are often organized by historical labels, mathematical assumptions, or the metaphor used to introduce an algorithm. These views remain useful for locating a method in the literature, but they can obscure structural relationships. Two algorithms with different names may preserve the same kind of search state, use similar information, and generate proposals through closely related operators; conversely, algorithms grouped under one broad family may differ substantially in memory, adaptation, selection, or model structure. We therefore use the algorithm census as an identity layer and a separate mechanism fingerprint as the unit of structural comparison.

The taxonomy is organized around three questions: what state is carried between evaluations, what information is used to update that state, and how is a new candidate or decision generated? Secondary axes describe memory, adaptation, selection, diversity control, constraint handling, derivative order, model type, parallelism, and dominant computational overhead. This representation does not assert that these fields completely determine algorithm behavior. Rather, it provides a common descriptive basis on which algorithms from otherwise separate literatures can be compared.

\subsection{State, information, and proposal geometry}

The most immediate distinction concerns the state maintained by the optimizer. Point-based methods maintain one incumbent or iterate, possibly augmented by a limited history. Nelder--Mead evolves a simplex, while DIRECT maintains a geometric partition of the domain. Population methods retain multiple candidate solutions and use interactions among them. Distribution-based methods, including CMA-ES, natural evolution strategies, estimation-of-distribution algorithms, and the cross-entropy method, elevate a parameterized sampling distribution to a primary search state. Surrogate methods maintain observations together with an explicit predictive or density model. ADMM maintains coupled primal and dual state, while sample-average approximation retains a sampled scenario representation of an uncertain objective.

A second axis describes the information transformed into search moves. Gradient and proximal methods use local differential information, with proximal operators incorporating composite structure. L-BFGS augments gradients with a short curvature history, coordinate descent restricts updates to selected coordinates, and online gradient descent uses the gradient of the currently revealed loss. Derivative-free local search instead obtains direction from objective comparisons, simplex geometry, polling directions, interpolation models, or neighborhood structure.

Population methods construct direction from relations among evaluated candidates. Differential evolution uses pairwise population differences; PSO combines personal and social best information through velocity dynamics; genetic algorithms use selected parents followed by recombination and mutation; ant-colony methods use accumulated pheromone information to bias constructive decisions. These mechanisms should not be collapsed into a single population-heuristic category because the information converted into a proposal is materially different.

Distribution-based search provides another transition. CMA-ES updates a covariance-structured search distribution using ranked samples and evolution paths. Natural evolution strategies update distribution parameters through a natural-gradient view, whereas estimation-of-distribution algorithms and the cross-entropy method refit probabilistic models from selected or elite samples. Surrogate optimization replaces or supplements direct search geometry with a model of observed performance. Gaussian-process Bayesian optimization uses posterior-dependent acquisition, TPE uses conditional density models, SMAC uses a model suited to mixed configuration spaces, and TuRBO localizes Bayesian optimization through trust regions. Hyperband is structurally different because its principal mechanism is resource allocation and successive elimination; BOHB combines resource allocation with model-based sampling.

\subsection{Selection, memory, and adaptation}

Selection determines which observations influence future search. Greedy replacement in differential evolution differs from probabilistic acceptance in simulated annealing, Pareto and diversity selection in multiobjective evolutionary algorithms, indicator selection in SMS-EMOA, and acquisition-driven selection in Bayesian optimization. Feasibility rules alter the ordering itself by incorporating constraint violation.

Memory ranges from an incumbent or short gradient history to archives, personal-best records, tabu lists, pheromone matrices, parameter memories, and learned probabilistic models. Tabu memory suppresses recently visited moves and supports diversification; personal-best memory in PSO influences velocity; the external archive in JADE changes the pool from which differential information can be drawn; SHADE stores successful control-parameter information; Bayesian optimization stores observations to update a surrogate.

Adaptation is therefore described by what is adapted and from which evidence. The DE lineage makes this distinction visible. Canonical DE uses externally selected control parameters. jDE self-adapts control parameters with individuals. SaDE uses accumulated search experience to adapt strategy choice and parameters. JADE introduces success-based parameter adaptation together with a current-to-p-best strategy and optional archive. SHADE adds an explicit history of successful parameter settings, and L-SHADE further introduces scheduled population-size reduction. These methods share differential variation while modifying the memory and control loops surrounding it.

\subsection{Constraints, uncertainty, and temporal structure}

Constraint handling is not a single auxiliary operation. COBYLA constructs local linear approximations of the objective and constraints inside a trust-region process. Feasibility rules alter dominance between candidates without requiring a problem-specific penalty coefficient. ADMM exposes constraints through operator splitting and coupled primal--dual updates. Robust optimization embeds uncertainty sets into the decision model, whereas sample-average approximation replaces an expectation by a finite scenario average. These mechanisms address different information structures and should not be compared solely under a generic constrained-optimization label.

Temporal information creates another distinction. Online gradient descent receives losses sequentially and is naturally evaluated through regret rather than only terminal objective value. Dynamic optimization methods face an objective or environment that changes over time. The census retains such research families even when the representative evidence layer does not promote every named dynamic method to a distinct fingerprint. Census coverage and mechanism evidence are therefore deliberately separate.

\subsection{Multiobjective selection as a structural layer}

Multiobjective optimization demonstrates that proposal generation alone is insufficient to characterize an algorithm. NSGA-II uses nondominated sorting and crowding to combine convergence and diversity. SPEA2 uses strength-based fitness, density estimation, and archive truncation. MOEA/D converts a multiobjective problem into linked scalar subproblems. NSGA-III introduces reference-point-based diversity management for many-objective settings, RVEA uses reference vectors, and SMS-EMOA uses hypervolume contribution as a selection signal. These algorithms may share conventional variation operators while differing principally in environmental selection and diversity geometry.

Accordingly, objective scope, selection, memory, and diversity control are encoded separately. This avoids treating all Pareto-based evolutionary methods as structurally interchangeable and permits mechanisms such as decomposition, reference directions, and indicators to be compared independently of the variation operator with which they are implemented.

\subsection{Genealogy versus structural similarity}

Genealogy and structural similarity answer different questions. Genealogy records documented historical or methodological descent: jDE, SaDE, JADE, SHADE, and L-SHADE belong to the broader DE development line, while BOHB explicitly combines model-based optimization with Hyperband-style resource allocation. Structural similarity instead asks whether normalized mechanism fields are close, regardless of documented descent. High structural similarity therefore does not establish ancestry, equivalence, or equivalent empirical behavior.

This distinction is particularly important for metaphor-driven optimization. The census retains a redundancy-audit lane containing methods introduced through different natural or social metaphors, but the evidence framework does not assume that each name defines a distinct mechanism. Promotion to a separate mechanism fingerprint requires an inspectable structural distinction supported by evidence. This rule preserves literature coverage without equating naming proliferation with mechanistic novelty.

\subsection{A mechanism-loop representation}

The fingerprints can be summarized through the descriptive update
\begin{equation}
\mathcal{S}_{t+1}=\mathcal{M}\!\left(\mathcal{S}_{t},\mathcal{I}_{t},\mathcal{P}(\mathcal{S}_{t},\mathcal{I}_{t};\theta_t),\mathcal{E}_{t}\right),
\label{eq:mechanism-loop}
\end{equation}
where \(\mathcal{S}_{t}\) is persistent search state, \(\mathcal{I}_{t}\) is information available at decision time \(t\), \(\mathcal{P}\) is the proposal mechanism, \(\theta_t\) is its control state, \(\mathcal{E}_{t}\) is newly observed objective, constraint, gradient, or model evidence, and \(\mathcal{M}\) is the memory, selection, and adaptation update. This notation is descriptive rather than a convergence theorem; it exposes which component changes when named algorithms are compared.

Under this view, algorithm development can change the state representation, information source, proposal geometry, evidence-retention rule, selection criterion, adaptation policy, or constraint and uncertainty model. A claimed new optimizer can therefore be examined at the component level. The same representation links the taxonomy to the evidence analysis that follows: empirical limitations can be associated with mechanism components without assuming that finite benchmark evidence establishes an intrinsic or universal failure.
